\scanpage{80}
\mspage{4}

est équivalente à celle des couples $\{t,i\}$, avec $t\in U_{0,3}(k)=\mathbb P^1_k(k)\setminus\{0,1,\infty\}$, et $i\in\mu_4^{\unsure{*}}$.

\deleted{On peut définir sur $\Sigma_J$ le fibré}

\[
\deleted{O_{\Sigma_J}(1)^{\natural}\deleted{\unsure{(}}=O_{\Sigma_J}(1)\wedge_{\mathord{\mathrm{I}\mkern-7mu\mu}_4}}
\]

Je repasse à une base quelconque, mais je travaille en car. rés. $\ne 2$ i.e. sur la base absolue \inserted{$S_0=$}$\operatorname{Spec}(\mathbb Z[1/2])$. Sur \deleted{$\unsure{\mathbb P^1_S|\Sigma_J|\Sigma_S}$} je définis l'opération $\unsure{\natural|\theta}$ $\natural$, sur les \deleted{modules} faisceaux de Modules, \deleted{\unsure{par|pour}} sur un schéma de base $S$

\[
M^\natural=M\unsure{\wedge|\otimes}_{\mathord{\mathrm{I}\mkern-7mu\mu}_{4S}}\underbrace{Q_{0S}}_{\deleted{\unsure{\mu_4|S|\text{un}}}}
\tag*{(22)}
\]

\margin{\uline{NB}\quad $M^\natural\simeq M\otimes_{O_S}\underline{O}_S^\natural$}

où ici on considère $Q_0$ comme un $\mathord{\mathrm{I}\mkern-7mu\mu}_4$-torseur sur la base absolue $S_0$.

On aura donc

\[
M^\natural\otimes M^\natural\simeq M^\flat\overset{\mathrm{def}}{=}M\wedge_{\mathord{\mathrm{I}\mkern-7mu\mu}_{2S}=\{\pm1\}_S}\mathord{\mathrm{I}\mkern-7mu\mu}_{4S}^*
\tag*{(23)}
\]

où \deleted{$\mathord{\mathrm{I}\mkern-7mu\mu}_4^*=\unsure{\mu}$}

\[
\mathord{\mathrm{I}\mkern-7mu\mu}_4^*=\mathord{\mathrm{I}\mkern-7mu\mu}_4\setminus\mathord{\mathrm{I}\mkern-7mu\mu}_2
\qquad\begin{aligned}
&\text{est considéré comme}\\[-2pt]
&\text{torseur sous }\mathord{\mathrm{I}\mkern-7mu\mu}_2,\ \text{i.e.}\\
&\text{rev. étale quadratique}\\
&\text{de la base absolue }S_0
\end{aligned}
\tag*{(24)}
\]

On s'intéresse donc, \deleted{\unsure{sur|au}} pour un revêtement
% archive sheet number: 79
